\documentclass[12pt,a4paper]{article}

\usepackage[utf8]{inputenc}
\usepackage[T2A]{fontenc}
\usepackage[russian]{babel}
\usepackage{amsmath,amssymb}
\usepackage{geometry}
\usepackage{graphicx}
\usepackage{float}

\geometry{margin=2.2cm}

\newcommand{\plotplaceholder}[2][0.9\textwidth]{%
    \IfFileExists{#2}{%
        \includegraphics[width=#1]{#2}%
    }{%
        \fbox{%
            \begin{minipage}[c][0.22\textheight][c]{#1}
            \centering
            Заглушка для графика\\[4pt]
            \texttt{\detokenize{#2}}
            \end{minipage}%
        }%
    }%
}

\begin{document}

\title{Оценка максимального правдоподобия}
\author{}
\date{}
\maketitle

\section*{Параметрическая модель}

В предыдущих разделах плотность оценивалась непараметрически: ядерной
оценкой или гистограммой. Теперь используем параметрический подход. Будем
считать, что форма распределения известна заранее: это равномерное
распределение, но его границы неизвестны.

Рассмотрим семейство плотностей
\[
f(x;A,B)=
\frac{1}{B-A}\mathbf 1_{\{A\le x\le B\}},
\qquad A<B.
\]
Истинная плотность из нашей задачи получается при
\[
A=-a,\qquad B=a,\qquad a=\sqrt{3}.
\]
Задача метода максимального правдоподобия состоит в следующем. По выборке
\[
X_1,\ldots,X_n
\]
сначала оцениваются параметры $A$ и $B$, а затем эти оценки подставляются в
параметрическую плотность.

\section*{ММП-оценки параметров}

Функция правдоподобия равна
\[
L(A,B)
=
\prod_{j=1}^n f(X_j;A,B)
=
\frac{1}{(B-A)^n}
\mathbf 1_{\{A\le X_{(1)},\, X_{(n)}\le B\}},
\]
где
\[
X_{(1)}=\min_{1\le j\le n}X_j,
\qquad
X_{(n)}=\max_{1\le j\le n}X_j.
\]
Чтобы правдоподобие было ненулевым, интервал $[A,B]$ должен содержать всю
выборку. При этом величина $(B-A)^{-n}$ тем больше, чем меньше длина
интервала. Поэтому максимальное правдоподобие достигается на самом коротком
интервале, содержащем выборку:
\[
\widehat A=X_{(1)},\qquad
\widehat B=X_{(n)}.
\]
Следовательно, ММП-оценка плотности имеет вид
\[
\widehat f^{ML}(x)
=
\frac{1}{X_{(n)}-X_{(1)}}
\mathbf 1_{\{X_{(1)}\le x\le X_{(n)}\}}.
\]
Это прямоугольная плотность на отрезке от минимального до максимального
наблюдения.

\section*{Интегральная квадратическая ошибка}

Обозначим случайный размах выборки через
\[
R=X_{(n)}-X_{(1)}.
\]
Поскольку все наблюдения лежат внутри истинного носителя, отрезок
\[
[X_{(1)},X_{(n)}]\subset[-a,a].
\]
Поэтому ошибка состоит из
двух частей: на отрезке $[X_{(1)},X_{(n)}]$ отличаются высоты двух
прямоугольников, а на оставшейся части истинного носителя оценка равна
нулю.

При фиксированной выборке
\[
\begin{aligned}
\int_{-\infty}^{\infty}
\bigl(\widehat f^{ML}(x)-f(x)\bigr)^2\,dx
&=
R\left(\frac{1}{R}-\frac{1}{2a}\right)^2
+(2a-R)\left(\frac{1}{2a}\right)^2\\
&=
\frac{1}{R}-\frac{1}{2a}.
\end{aligned}
\]
Теперь нужно взять математическое ожидание по выборке.

\section*{Распределение размаха}

Перейдём к нормированному размаху
\[
W=\frac{R}{2a}.
\]
Для выборки из равномерного распределения на отрезке плотность $W$ равна
\[
p_W(w)=n(n-1)w^{n-2}(1-w),
\qquad 0<w<1.
\]
Тогда при $n>2$
\[
\begin{aligned}
\mathbb E\frac{1}{W}
&=
n(n-1)\int_0^1 w^{n-3}(1-w)\,dw\\
&=
n(n-1)
\left(\frac{1}{n-2}-\frac{1}{n-1}\right)
=
\frac{n}{n-2}.
\end{aligned}
\]
Так как $R=2aW$,
\[
\mathbb E\frac{1}{R}
=
\frac{1}{2a}\mathbb E\frac{1}{W}
=
\frac{1}{2a}\cdot\frac{n}{n-2}.
\]
Значит,
\[
\operatorname{MISE}_{ML}
=
\mathbb E\int(\widehat f^{ML}(x)-f(x))^2\,dx
=
\frac{1}{2a}\left(\frac{n}{n-2}-1\right)
=
\frac{1}{a(n-2)}.
\]
При $n\le 2$ величина $\mathbb E(1/R)$ бесконечна, поэтому средняя
интегральная квадратическая ошибка не конечна.

\section*{Нормированная ОИСКО}

Нормируем ошибку на $\|f\|_2^2$. Для равномерной плотности
\[
\|f\|_2^2
=
\int_{-a}^{a}\frac{1}{4a^2}\,dx
=
\frac{1}{2a}.
\]
Следовательно, при $n>2$
\[
\operatorname{RMISE}_{ML}(n)
=
\frac{\operatorname{MISE}_{ML}}{\|f\|_2^2}
=
\frac{2}{n-2}.
\]
Эта формула показывает, что ошибка ММП-оценки убывает как величина порядка
$1/n$.

\section*{Графики}

На первом графике показана точная формула
$\operatorname{RMISE}_{ML}(n)=2/(n-2)$.

\begin{figure}[H]
\centering
\plotplaceholder[0.82\textwidth]{mle_rmise_vs_n.png}
\caption{Нормированная ОИСКО ММП-оценки в зависимости от объёма выборки.}
\end{figure}

Следующий график показывает, как выглядит сама ММП-оценка плотности. Синим
показана оценка на отрезке от минимума до максимума выборки, красным ---
истинная плотность.

\begin{figure}[H]
\centering
\plotplaceholder[0.92\textwidth]{mle_density_examples.png}
\caption{Примеры ММП-оценки плотности для разных объёмов выборки.}
\end{figure}

Последний график является численной проверкой аналитической формулы:
точная кривая сравнивается со средним значением нормированной ошибки,
посчитанным по большому числу случайных выборок.

\begin{figure}[H]
\centering
\plotplaceholder[0.82\textwidth]{mle_mc_check.png}
\caption{Сравнение аналитической формулы с численной проверкой.}
\end{figure}

\section*{Сравнение с предыдущими оценками}

Теперь сравним ММП-оценку с ранее рассмотренными непараметрическими
оценками. Во всех случаях используется одна и та же нормированная ОИСКО,
поэтому значения можно сравнивать напрямую.

\[
\begin{array}{c|c|c|c|c}
n
& \text{ядерная}
& \text{гист., }\Delta=0.25
& \text{ММП}
& \text{гист., }\Delta=0\\
\hline
30  & 0.1206 & 0.1284 & 0.0714 & 0\\
100 & 0.0725 & 0.0778 & 0.0204 & 0
\end{array}
\]

\begin{figure}[H]
\centering
\plotplaceholder[0.82\textwidth]{mle_methods_comparison.png}
\caption{Сравнение минимальной нормированной ОИСКО для разных методов оценки.}
\end{figure}

ММП-оценка лучше ядерной оценки и сдвинутой гистограммы, потому что она
использует правильную параметрическую модель: заранее известно, что
плотность равномерная, а неизвестны только границы носителя. Нулевые
значения для гистограммы при $\Delta=0$ являются специальным эффектом
идеально согласованной сетки и не являются типичным поведением гистограмм.

\end{document}
